\documentclass{article}

\usepackage[T1]{fontenc}

\usepackage{booktabs,graphicx,siunitx}

\begin{document}

\section{One-way ANOVA + post-hoc}

\begin{tabular}{llll}
\toprule
Group & N & Mean & Std. Dev. \\
\midrule
Blended & 80 & 65.49 & 6.85 \\
Flipped & 80 & 70.75 & 6.58 \\
Lecture & 80 & 60.46 & 7.57 \\
\bottomrule
\end{tabular}

\begin{tabular}{lllllll}
\toprule
Source & SS & df & MS & F & p & $\eta$\textsuperscript{2} [99\% CI] \\
\midrule
Between & 4234.40 & 2 & 2117.20 & 43.07 & $<$.001 & 0.27 [0.16, 1.00] \\
Within & 11650.16 & 237 & 49.16 &  &  &  \\
\bottomrule
\end{tabular}

\begin{tabular}{lllll}
\toprule
Pair & Mdiff & SE & padj & 99\% CI \\
\midrule
Blended - Flipped & $-$5.26 & 1.11 & $<$.001 & [$-$8.52, $-$2.00] \\
Blended - Lecture & 5.03 & 1.11 & $<$.001 & [1.77, 8.29] \\
Flipped - Lecture & 10.29 & 1.11 & $<$.001 & [7.03, 13.55] \\
\bottomrule
\end{tabular}

\includegraphics{figures/01_one-way-anova/figure_means-plot.png}

The F and its p tell you whether the groups differ overall; $\eta$\textsuperscript{2} (or $\omega$\textsuperscript{2}) is the effect size. If significant, the post-hoc table shows which specific pairs differ, with multiplicity-adjusted p-values. Your significance threshold ($\alpha$) is 0.05.

A one-way ANOVA gave F(2,237)=43.07, p $<$ .001, $\eta$\textsuperscript{2}=.27 [.16, 1.00]. Tukey HSD post-hoc tests showed\ldots{}

Statistical basis: F-test for equality of several means (ANOVA) (Fisher, R. A. (1925). Statistical Methods for Research Workers. Oliver \& Boyd.); Tukey HSD post-hoc comparison (Tukey, J. W. (1949). "Comparing individual means in the analysis of variance." Biometrics, 5(2), 99-114.); Bonferroni-corrected post-hoc comparison (a 'posthoc' choice) (Dunn, O. J. (1961). "Multiple comparisons among means." Journal of the American Statistical Association, 56(293), 52-64.); Scheff\'{e}-corrected post-hoc comparison (a 'posthoc' choice) (Scheff\'{e}, H. (1953). "A method for judging all contrasts in the analysis of variance." Biometrika, 40(1-2), 87-104.); Levene's test for equal variances (Brown-Forsythe median-centered variant) (Levene, H. (1960). "Robust tests for equality of variances." In I. Olkin (Ed.), Contributions to Probability and Statistics (pp. 278-292). Stanford University Press.); Brown-Forsythe median-centered variant of the homogeneity-of-variance test (Brown, M. B., Forsythe, A. B. (1974). "Robust tests for the equality of variances." Journal of the American Statistical Association, 69(346), 364-367.); $\eta$\textsuperscript{2} effect-size benchmarks (Cohen, J. (1988). Statistical Power Analysis for the Behavioral Sciences (2nd ed.). Routledge.).

\end{document}